%% Trottinette électrique en montée sur une pente de 14 % (soit 14 m pour 100 m à l'horizontale).
%% Isolement de l'ensemble {usager + trottinette} : poids P en G, action du sol en A (roue avant
%% motrice : traction + réaction normale) et en B (roue arrière : réaction normale).
%% L'angle alpha n'est pas coté en degrés : sa valeur est le résultat demandé à l'élève.
\documentclass[tikz,border=4pt]{standalone}
\usepackage{liaisons}
\begin{document}
\begin{tikzpicture}[x=1cm,y=1cm,
  vect/.style={-{Stealth[length=6pt]}, line width=1.2pt, solideA},
  proj/.style={-{Stealth[length=4pt]}, line width=0.7pt, solideA!70,
               dash pattern=on 2.5pt off 2pt},
  piece/.style={line width=1.5pt, solideB, line cap=round, line join=round},
  cote/.style={line width=0.6pt, solideD, {Stealth[length=4pt]}-{Stealth[length=4pt]}},
  fin/.style={line width=0.5pt, black!60, dash pattern=on 2.5pt off 2pt}]

\useasboundingbox (-0.6,-1.6) rectangle (8.0,3.7);
\def\al{7.97}                       % pente 14 % : tan(alpha) = 0,14
\coordinate (O) at (0.4,0);
\coordinate (B) at (2.55,0.301);    % contact roue arrière / sol
\coordinate (A) at (4.62,0.591);    % contact roue avant motrice / sol
\coordinate (G) at (3.60,2.30);     % centre de gravité de l'ensemble usager + trottinette

%% ---------------- la pente et le sol ------------------------------------
\draw[line width=1.2pt, solideE] (O) -- (7.6,1.008);
\foreach \s in {0.15,0.45,...,7.1} {
  \draw[line width=0.5pt, solideE!80] ($(O)+(\al:\s)$) -- ++(-0.16,-0.24);}
\draw[fin] (O) -- ++(2.5,0);
\draw[line width=0.6pt, solideE] ([shift=(0:1.85)]O) arc (0:\al:1.85);
\node[font=\small, text=solideE, anchor=west, inner sep=1pt] at ([shift=(3.6:1.92)]O) {$\alpha$};
\node[font=\small, text=solideE, anchor=north west] at (0.55,-0.30) {pente 14 \%};

%% ---------------- triangle de pente (100 m / 14 m) -----------------------
\draw[line width=0.7pt, solideD] (3.30,-1.05) -- (6.30,-1.05) -- (6.30,-0.63) -- cycle;
\node[font=\scriptsize, text=solideD, anchor=north] at (4.80,-1.08) {100 m};
\node[font=\scriptsize, text=solideD, anchor=west, inner sep=2pt] at (6.35,-0.84) {14 m};

%% ---------------- la trottinette -----------------------------------------
\draw[piece] (3.00,0.60) -- (4.20,0.78);                       % plateau
\draw[piece] (2.51,0.577) circle (0.28);                       % roue arrière
\draw[piece] (4.58,0.867) circle (0.28);                       % roue avant motrice
\draw[piece] (4.58,0.867) -- (5.18,2.05);                      % colonne de direction
\draw[piece] (4.93,1.90) -- (5.43,2.20);                       % guidon
%% silhouette de l'usager
\draw[line width=2.2pt, black!45, line cap=round, line join=round]
  (3.35,0.66) -- (3.52,1.55) -- (3.75,0.72)                    % jambes
  (3.52,1.55) -- (3.44,2.62)                                   % tronc
  (3.44,2.55) -- (4.30,2.30) -- (5.05,2.00);                   % bras vers le guidon
\draw[line width=1.2pt, black!45, fill=black!12] (3.40,2.83) circle (0.20);

%% ---------------- poids et projections ------------------------------------
\fill (G) circle (2pt);
\node[font=\small, anchor=east, inner sep=3pt] at (3.52,2.34) {$G$};
\draw[vect] (G) -- ++(0,-1.6) node[anchor=west, inner sep=3pt, font=\small, text=solideA] {$\vec{P}$};
\draw[proj] (G) -- ++(\al+180:0.222);
\draw[proj] (G) -- ++(\al-90:1.584);
\draw[line width=0.4pt, black!45, dash pattern=on 2pt off 2pt]
  ($(G)+(\al+180:0.222)$) -- ++(\al-90:1.584);
\node[font=\scriptsize, text=solideA, anchor=east, inner sep=2pt] at (3.30,2.42)
     {$-mg\sin\alpha\,\vec{u}$};
\node[font=\scriptsize, text=solideA, anchor=west, inner sep=3pt] at (4.02,1.55)
     {$-mg\cos\alpha\,\vec{v}$};

%% ---------------- actions du sol -------------------------------------------
\draw[vect] (A) -- ++(\al:1.15)
     node[anchor=north west, inner sep=1pt, font=\small, text=solideA] {$\vec{F}_{\text{tract}}$};
\draw[vect] (A) -- ++(\al+90:1.05)
     node[anchor=south west, inner sep=1pt, font=\small, text=solideA] {$\vec{R}_{\text{av}}$};
\draw[vect] (B) -- ++(\al+90:0.95)
     node[anchor=south east, inner sep=1pt, font=\small, text=solideA] {$\vec{R}_{\text{ar}}$};
\fill (A) circle (1.6pt); \node[font=\small\bfseries, anchor=north, inner sep=3pt] at (A) {$A$};
\fill (B) circle (1.6pt); \node[font=\small\bfseries, anchor=north, inner sep=3pt] at (B) {$B$};

%% ---------------- repère local et vitesse ----------------------------------
\begin{scope}[shift={(6.35,1.15)}]
  \draw[-{Stealth[length=5pt]}, line width=0.7pt, solideF] (0,0) -- (\al:0.8)
       node[anchor=north west, inner sep=1pt, font=\small, text=solideF] {$\vec{u}$};
  \draw[-{Stealth[length=5pt]}, line width=0.7pt, solideF] (0,0) -- (\al+90:0.8)
       node[anchor=south east, inner sep=1pt, font=\small, text=solideF] {$\vec{v}$};
\end{scope}
\draw[-{Stealth[length=6pt]}, line width=1pt, solideD] (5.95,2.55) -- ++(\al:1.15);
\node[font=\scriptsize, text=solideD, anchor=south east, inner sep=2pt] at (7.15,2.70)
     {$\vec{V}_{G,t/\text{sol}}$ constante};

%% ---------------- cote du diamètre de roue ----------------------------------
\draw[line width=0.5pt, solideD] (4.78,0.68) -- (5.55,0.30);
\node[font=\scriptsize, text=solideD, anchor=west, inner sep=1pt] at (5.55,0.28) {diamètre 216 mm};
\end{tikzpicture}
\end{document}
